\chapter[A decision before and after disclosure]{A decision before and after disclosure: an exact learning example}
\label{sim:learning-case}

Learning matters only through a specified relation between hidden facts, observations and decisions. This chapter constructs that relation explicitly. The example is finite, so every policy and probability can be checked by enumeration. It demonstrates why two companies with equal current statements can require different information states, and why preserving only unconditional future cash distributions is insufficient for adaptive decisions.

\section{Two hidden exposures}
A company must choose a hedge before a future settlement. Its exposure type is either A or B. Under type A, hedge A avoids a loss of 10 while hedge B avoids none. Under type B, hedge B avoids a loss of 10 while hedge A avoids none. Each hedge costs 2. Doing nothing costs nothing and leaves the loss 10.

For the first calculation, assume type A and B each have probability one half. This probability is a supplied teaching assumption. The company has enough cash to execute either hedge and can choose at most one. Current financial statements do not distinguish the types.

Net loss after the correct hedge is 2, after the wrong hedge is 12, and after no hedge is 10. Without further information, either hedge has expected loss 7, so a risk-neutral expected-loss policy chooses either hedge rather than doing nothing. Symmetry gives two equally optimal actions. The model should not invent a preference between them.

\section{A noisy signal}
Before choosing, the actor can receive a signal labeled A or B. Let the signal be correct with probability $q$, symmetrically across types, with $q\geq1/2$. Bayes' rule with equal priors gives posterior probability $q$ for the signaled type. Following the signal yields expected loss
\begin{equation}
 L_{\rm signal}(q)=2+10(1-q)=12-10q.
 \label{sim:signal-loss}
\end{equation}
The first term is hedge cost. The second is the loss when the signal is wrong. At $q=0.9$, expected loss is 3; at $q=0.5$, it is 7, matching no informative signal.

The gross value of the signal relative to the best uninformed hedge is $7-(12-10q)=10q-5$. A signal costing $k$ is worthwhile under this criterion when $k<10q-5$. At $q=0.9$, the maximum willingness to pay is 4. This calculation requires the signal to arrive before the action and the actor to be able to respond.

\section{A complete joint probability table}
For $q=0.9$, the four type-signal combinations have probabilities
\begin{center}
\begin{tabular}{llrr}
\toprule Type&Signal&Probability&Loss following signal\\\midrule
A&A&0.45&2\\
A&B&0.05&12\\
B&A&0.05&12\\
B&B&0.45&2\\\bottomrule
\end{tabular}
\end{center}
The expected loss is $0.9(2)+0.1(12)=3$. The signal is marginally A or B with equal probability. Type is also marginally A or B with equal probability. The information lies in their dependence.

If a compression retains these two marginals but makes them independent, it changes each joint cell to 0.25. Following the signal then has expected loss 7. The compressed model has destroyed the signal's information while preserving both marginal distributions. This is why the joint successor-observation law appears in the exact aggregation condition.

\section{Two observation technologies with equal current reports}
Consider company state U, which will receive the informative signal with $q=0.9$, and state V, which will receive an uninformative signal with $q=0.5$. Their current cash, debt and income are identical. Their hidden type priors are identical. Before signal arrival, a fixed hedge has the same expected loss 7 in both.

An adaptive policy following the signal has expected loss 3 in U and 7 in V. Merging U and V because current reports and fixed-action values agree loses a decision-relevant difference. The future observation technology is part of state for this question.

This finite example gives an accounting-oriented counterpart to the policy-versus-action-sequence distinction in \citet[Section 3.1]{givan2003}. The numerical construction is local. The source supplies the general warning that policies react to intermediate information and therefore require a stronger equivalence than fixed sequences.

\section{What happens when the signal arrives late}
Suppose the hedge must be committed at date 1 and the signal arrives at date 2. The informative technology no longer improves that hedge decision. A policy that follows the date-2 signal but executes the hedge at date 1 is infeasible.

The expected-loss value returns to 7 for the best uninformed hedge, even though the analyst eventually learns much more about the type. Information has retrospective explanatory value but no value for this already committed action. A simulation must represent action deadlines and settlement lead times, not only observation dates.

If the contract permits cancellation or adjustment at date 2 for cost $c$, some signal value can return. The calculation then compares adjustment benefits with $c$ and available liquidity. The ability to revise is a contract term or policy permission, not a default property of every action.

\section{A cancellation calculation}
Suppose hedge A is chosen initially. After receiving the signal, the actor may switch to B at extra cost $c$. With signal B, posterior probability of type B is $q$. Keeping A has conditional expected loss $2+10q$. Switching has conditional expected loss $2+c+10(1-q)$.

Switching is preferred when
\begin{equation}
 2+c+10(1-q)<2+10q\quad\Longleftrightarrow\quad c<20q-10.
\end{equation}
At $q=0.9$, the threshold is 8. The switching opportunity is exercised only after signal B under the symmetric setup. Its unconditional value must also account for the probability of that signal.

If the initial hedge cannot be unwound and the second hedge creates an offsetting or additional exposure rather than replacing it, the formula changes. A model should not assume that a new action erases an old contract. This is the same event-persistence issue seen in debt amendments.

\section{Unknown signal quality}
Now suppose evidence supports only $q\in[0.6,0.9]$. The expected loss from following the signal lies in $[3,6]$ under the same prior and symmetry assumptions. The gross information value lies in $[1,4]$. A signal priced at 0.5 improves expected loss for every admitted $q$; one priced at 3 is worthwhile only for some qualities.

These are bounds over a parameterized probabilistic family. They are not probabilities over signal quality. Assigning a uniform distribution to $q$ would add another assumption. A robust decision can use the worst admitted quality, while a Bayesian decision can use a prior over it if justified.

If the quality depends on exposure type, one scalar $q$ is insufficient. Let sensitivity and specificity differ. Posterior probabilities then depend on the prior and both error rates. A report that is usually correct overall can still be misleading for a rare type. The observation model should preserve the relevant asymmetry.

\section{Unequal prior probabilities}
Let prior probability of A be $p$ and signal accuracy remain symmetric at $q$. After signal A, the posterior is
\begin{equation}
 p_A^+=\frac{pq}{pq+(1-p)(1-q)}.
\end{equation}
After signal B, it is
\begin{equation}
 p_B^+=\frac{p(1-q)}{p(1-q)+(1-p)q}.
\end{equation}
The actor chooses hedge A when the posterior probability of A exceeds one half, hedge B when below, and is indifferent at equality under the symmetric loss structure.

If $p$ is very high, even signal B may not overturn the prior. Following the signal mechanically is then suboptimal. For example, $p=0.95$ and $q=0.9$ gives posterior A after signal B equal to $0.095/(0.095+0.045)$, about 0.6786. Hedge A remains preferred.

This illustrates why a signal label is not itself a decision. The policy combines the observation with prior information and payoffs. In an unweighted reconstruction, there may be no justified prior at all, so a posterior calculation would be inappropriate until that modeling choice is admitted.

\section{An unweighted version}
Suppose both types remain possible and no probabilities are assigned. An exact signal identifying the type allows the correct hedge in every world, producing loss 2. Without a signal, either hedge has worst-case loss 12, while doing nothing has worst-case loss 10. A minimax-loss policy therefore chooses no hedge before information and the correct hedge after exact information.

The criterion differs from the earlier expected-loss decision, which chose a hedge without a signal. This is not a contradiction. The uncertainty representation and decision rule changed. A system should report the dependence rather than label one answer universally optimal.

If a noisy signal can be wrong in either type and only support is known, both types may remain possible after either signal. The worst-case benefit can be zero. A probability of correctness above one half is a probabilistic restriction; it cannot be inferred from support alone.

\section{Management and public investors}
Suppose management already knows the type from internal contracts, while public investors see only aggregate statements. Management can choose the correct hedge with loss 2 even before public disclosure. The analyst should simulate that informed behavior if it is the declared policy.

A public announcement can then affect investor beliefs and financing offers without changing management's hedge choice. The same information event has different economic consequences for different actors. A single global belief state cannot represent this distinction.

If management's action itself is observable, investors may learn from it. A hedge choice can signal the hidden type under an assumed policy. That inference depends on the policy model. If management sometimes hedges for other reasons, the action is not a perfect signal. Endogenous disclosure and signaling require a richer model than an exogenous likelihood.

\section{Endogenous observations}
A company can choose when and what to disclose within its obligations. A lender can request information before making an offer. An auditor can discover discrepancies with probability depending on effort. These observations depend on actions and hidden state.

The observation kernel should then be written conditional on action as well as state. Treating disclosure as independent can distort the value of monitoring or the selection of companies that reveal information. A missing report may itself be informative only under an explicit model of why reports are missing.

AccDSL need not assume a strategic disclosure game in every application. It should allow exogenous observation models as a bounded profile and state when they are inadequate. If the query concerns financing responses to disclosure, the endogenous channel may be material and should remain an open modeling requirement.

\section{Persistent type and repeated signals}
Suppose two conditionally independent signals of accuracy 0.9 arrive about the same fixed type with equal prior. Two A signals give posterior A
\begin{equation}
 \frac{0.9^2}{0.9^2+0.1^2}=\frac{81}{82}.
\end{equation}
An A followed by B returns posterior one half under the symmetry and independence assumptions. These updates concern one persistent type, not a newly drawn type for each signal.

If the signals are copies of the same underlying source, conditional independence is false. Counting a duplicated report as independent evidence would overstate confidence. Two identical copies should not square the likelihood unless the observation process actually supplies independent information.

Source provenance therefore matters statistically. A news summary quoting a filing is not necessarily another independent measurement of the loan type. The system should preserve derivation links so that a likelihood model can avoid double counting.

\section{Correlated errors}
Suppose both signals share a common extraction error. With probability 0.1 the type label is reversed in the source, and both downstream reports repeat it. Two matching reports then have accuracy 0.9, not $81/82$. The agreement supplies no additional independent confirmation.

A model with separate independent error terms would produce unjustified confidence. The remedy is a shared latent error or a source-dependence relation. If the dependence is unknown, sensitivity over plausible observation models is more defensible than assuming independence because it is convenient.

The same issue arises when several financial databases reproduce one filing or when multiple ratios use the same extracted amount. More records do not necessarily mean more evidence. The model should count informational content through its observation relation, not file count.

\section{A belief-state calculation}
In the finite probabilistic example, the posterior probability $p_t$ of type A summarizes the evidence relevant to the hedge decision, provided the signal model and loss structure remain fixed. Expected losses are $2+10(1-p_t)$ for hedge A, $2+10p_t$ for hedge B and 10 for no hedge.

The policy chooses the smallest of these expressions. Hedge A is preferred to B when $p_t\geq1/2$. It is preferred to no hedge when $p_t\geq0.2$. Hedge B is preferred to no hedge when $p_t\leq0.8$. Thus at least one hedge is preferred throughout $[0,1]$ in this particular cost structure, with the choice switching at one half.

This compact belief is sufficient because the future depends on the past only through the fixed type posterior under the declared model. If signal quality is itself unknown or past actions create commitments, $p_t$ alone may cease to be sufficient. The state must then include those additional quantities or their joint belief.

\section{Compression tests from the example}
A proposed compression should preserve the joint type-signal distribution, action availability, costs and timing. It should reproduce expected loss 3 for informative signals and 7 for uninformative signals. It should refuse to merge those technologies for an adaptive policy query.

A separate test can merge states differing only in an irrelevant customer label when neither payoffs nor observations depend on that label. This establishes a valid reduction within the small model. Positive and negative examples are both needed: a system that never merges is safe from one error but does not solve the computational problem.

Changing the query to one fixed hedge can make a coarser compression sufficient. The approval should be scoped to that query. Reusing it for information acquisition or adaptive hedging without rechecking would repeat the original mistake.

\section{What the case teaches about trees}
The full tree has hidden types, signals and actions. Hidden type branches need not be visible to the actor. Signal branches define information partitions, and actions before a signal must be shared across indistinguishable worlds. A drawing that places separate decisions under each hidden type can accidentally imply forbidden knowledge.

A belief-state or symbolic implementation can encode the same problem without storing every path. It must still preserve the information partition. The tree is a useful explanatory projection of that structure, not its only possible storage form.

The central requirement is therefore not one point per possibility. It is preservation of the joint relation among persistent facts, observations and feasible decisions. This example gives exact numbers against which a proposed representation can be tested before it is applied to a much larger company.
